\section{结论}

本文从一个具体的制度反差出发：中国对烟花爆竹同时握有一个三十年未调整的从价消费税，和一套逐城逐年调整的禁放限放体系。为什么后者被不断叠加，而前者从未被重新校准？

三点发现。

\textbf{第一，外部性的形态决定了工具的适用性。} 正月初一 \pmm{2.5} 相对基线平均上升 43.1 \ugm（平衡面板），中位数 28.1，城际差异从负值到数百不等。方差分解显示，城市间差异只解释 22.0\%，年份间只解释 12.7\%，而 65.3\% 的变异落在城市 $\times$ 年份这一维度上——恰好是一个全国统一、按年适用的税率在构造上够不到的那一维。

\textbf{第二，这次因果识别失败了，而失败的原因是结构性的。} 可用作识别的市级切换城市只有 7 个，且 93\% 的城市-年没有政策记录；安慰剂结果变量显著；估计量对燃放日没有亲和性；禁令总是与同期的工业减排同文下达。\textbf{报告这次失败比掩盖它更有价值}，因为它说明这类政策在现行制度安排下无法被单独评估。

\textbf{第三，价格工具的问题不在税率而在解析度。} 全国统一、按年、生产环节的从价税，在城市与时间两个维度上的分辨率都是零。因此“提高税率”不是本文能支持的处方——它改变价格水平，不改变工具的解析度。庇古校准也给不出可引用的比值（跨约三个数量级）；它的失败源于两个没有来源的参数，因此不能单独证明异质性——但与第~\ref{sec:variance}~节的方差分解合起来，它指向同一件事：对这个源，“每公斤的边际损害”无法被一个全国统一的税率表达。

如果要把这三点压缩成一句：\textbf{当损害是局部的、短暂的、分散在千万个家庭除夕夜里的，一个统一的税率既收不到它，也算不出它。} 而地方政府选择的禁令，其成本恰恰藏在那些不被核算的警力工时里——\textbf{成本不进价格，效果也不单独核算。} 这两件事是同一件事。
